\scanpage{122}
\uline{Corollaire 3} Dans les conditions de la proposition, \deleted{considérons le groupe des $D(X)$ des diviseurs loc. principaux} \inserted{Soit $D$ un diviseur loc.} principal sur $X$, \deleted{\unsure{supposons les}} \inserted{soit $n_i$ l'entier déf par $D\cap\underline{\mathcal O}_{y_i}=\unsure{\sigma}(f_{y_i})^{n_i}$ i.e. le coefficient de $Y_i$ dans $D$.} \inserted{Supposons les} conditions suivantes satisfaites
\[
\text{(i)}\quad D\cap X_f=0
\]
(ii) \deleted{$D\cap\underline{\mathcal O}_{y_i}=\unsure{\sigma}(f_i)$} Tous les $n_i$ égaux \inserted{à un \unsure{même} $n$}

Alors $D=n\deleted{(f)}.\operatorname{div}(f)$

\uline{Corollaire 5} (\uline{Non démontré}) \inserted{Dans les conditions de} \deleted{\uline{Prop 4}} Prop, supposons \inserted{(ii) Pour tout $X_\alpha$, $Y\cap X_\alpha$ est irréd. (iii)} $Y$ connexe par la ``relation \deleted{d'équivalence} \inserted{d'incidence}'' suivante: $Y_i$ et $Y_j$ sont \uline{incidents} s'il existe une composante irréductible \inserted{$Z_{ij}\unsure{\subset}Y_i\cap Y_j$} \inserted{\unsure{contenant un|contenue dans un} $Z_{ij}$ de $\unsure{Y_i|X_i}\cap\unsure{Y_j|X_j}$} \deleted{\unsure{qui rencontre} $X_f$. \unsure{telle que la partie}}

Alors \deleted{tout} diviseur \deleted{\unsure{sur}} $D$ sur $X$ tel que $D\cap X_{\unsure{f|g}}=0$ est de la forme $n\operatorname{div}(f)$.

\margin{\uline{N.B} les hypothèses \unsure{faites} impliquent certainement que \unsure{l'on a une correspondance} \uline{biunivoque} \unsure{des composantes de} $X$ \unsure{aux} $Y_i$.
\diagram{\begin{tikzcd}X_\alpha \arrow[r,mapsto] & X_\alpha\cap Y\end{tikzcd}}
\unsure{1°)} $X_\alpha\cap Y$ est un $Y_i$;
\unsure{2°)} $X_\alpha\cap Y=X_\beta\cap Y\Rightarrow\alpha=\beta$
}

\inserted{Cas particulier. Dans les conditions de la Prop}
\uline{Cor \unsure{4|5}} Supposons \deleted{$X$ \unsure{irréd.}} $Y$ irréductible (donc $X$ \unsure{irréd.}
\deleted{Alors} tout $D$ tel que $D\cap X_f=0$, Alors $D=n\operatorname{div}(f)$
\margin{121}
